\documentclass[UTF8,a4paper,11pt]{ctexart}
\usepackage[a4paper,margin=2.3cm]{geometry}
\usepackage{amsmath,amssymb,mathtools}
\usepackage{enumitem}
\usepackage[hidelinks]{hyperref}
\usepackage{fancyhdr}
\setlist[itemize]{leftmargin=2.2em,itemsep=0.25em,topsep=0.35em}
\setlength{\parindent}{2em}\setlength{\parskip}{0.35em}
\pagestyle{fancy}\fancyhf{}\fancyhead[L]{离散数学与结构（I）}\fancyhead[R]{Lec 13}\fancyfoot[C]{\thepage}\setlength{\headheight}{14pt}
\title{生成函数}\author{zhangjunshuo}\date{2026 年 10 月}
\begin{document}
\maketitle
\begin{quote}
本讲根据 2026 年 10 月 27 日的两段课堂转写、两页手写原稿和课程教材整理。手写稿是公式清单（基本生成函数表、运算、找零、Catalan、$3\times n$ 覆盖的辅助状态、指数生成函数、Bernoulli 数），转写补充口述解释、图形分解与幂和推导；整理版按"取并集"处理，两边冲突或疑似笔误处给出勘误。

本讲主线：普通生成函数（定义/运算/多米诺/找零/部分分式）$\to$ 组合方程（Catalan、$3\times n$ 覆盖）$\to$ 指数生成函数（二项卷积、$\exp$/集合、教室与标号树）$\to$ 幂和与 Bernoulli 数。

每个分节开头都先交代本节用到的记号；所有符号在首次使用处给出定义。
\end{quote}
\tableofcontents\clearpage

\section{普通生成函数：定义、基本例子与运算}


生成函数把\textbf{数列}变成\textbf{形式幂级数}，从而把"数数列"的问题翻译成"做代数/微积分"的问题。这一节建立普通生成函数（ordinary generating function，OGF）的定义、基本例子与运算表，并用多米诺覆盖与找零两个问题说明它的用法。

\textbf{本节用到的记号：} $[x^n]A(x)$ 表示形式幂级数 $A(x)$ 中 $x^n$ 的系数；$\delta_{n,0}$ 在 $n=0$ 时为 $1$、否则为 $0$；$\binom nk$ 为二项式系数；$n^{\underline{k}}$ 为下降阶乘。若无特别说明，$\sum$ 的下标都从 $0$ 开始。

\subsection{定义：数列与形式幂级数}


\textbf{定义（普通生成函数）。} 给定数列 $\{a_n\}_{n\ge0}$，称形式幂级数

\[
A(x)=a_0+a_1x+a_2x^2+\cdots=\sum_{n\ge0}a_nx^n
\]

为它的\textbf{普通生成函数}。反过来，已知 $A(x)$ 时，$a_n=[x^n]A(x)$。

这里的"形式"指：我们只把 $x$ 当成一个记号，按幂级数的加法、乘法、求导、复合等代数规则运算，\textbf{不关心它是否收敛}。若某个具体问题中 $x$ 取实数、级数收敛，那当然更好（可以把 $A(x)$ 当成函数做分析），但组合计数里的恒等式在形式上就成立，不需要收敛性。这一点在后面的"有理函数"与"指数生成函数"里同样适用。

\textbf{例.} 常数数列 $1,1,1,\dots$ 的生成函数是

\[
1+x+x^2+\cdots=\frac{1}{1-x}.
\]

严格说，这是形式幂级数的等式：$(1-x)(1+x+x^2+\cdots)=1$，两边按幂级数乘法计算即可，不需要 $|x|<1$。

\subsection{一些最基本的生成函数}


下面这张表要熟记（都可以由等比级数或二项式定理直接得到）：

\begin{itemize}
\item $\{1,0,0,0,\dots\}\longleftrightarrow 1$；
\item $\{0,\dots,0,1,0,\dots\}$（第 $m$ 位为 $1$）$\longleftrightarrow x^m$；
\item $\{1,1,1,\dots\}\longleftrightarrow \dfrac{1}{1-x}$；
\item $\{1,\alpha,\alpha^2,\dots\}\longleftrightarrow \dfrac{1}{1-\alpha x}$；
\item $\{1,0,1,0,1,\dots\}\longleftrightarrow \dfrac{1}{1-x^2}$；
\item $\{1,0,\dots,0,1,0,\dots\}$（每隔 $p$ 位一个 $1$）$\longleftrightarrow \dfrac{1}{1-x^p}$；
\item $\left\{\binom nk\right\}_{n\ge0}\longleftrightarrow \dfrac{x^k}{(1-x)^{k+1}}$；
\item $\left\{\binom{n+k-1}{k-1}\right\}_{n\ge0}\longleftrightarrow \dfrac{1}{(1-x)^k}$；
\item $\left\{\binom{n}{k}\alpha^{n-k}\right\}_{n\ge0}\longleftrightarrow \dfrac{x^k}{(1-\alpha x)^{k+1}}$。
\end{itemize}


\textbf{推导示例.} 最后两行来自

\[
\frac{1}{(1-x)^k}=(1+x+x^2+\cdots)^k=\sum_{n\ge0}\binom{n+k-1}{k-1}x^n,
\]

因为 $x_1+\cdots+x_k=n$ 的非负整数解个数是 $\binom{n+k-1}{k-1}$（第 4 节"球与桶"的隔板法）。而

\[
\frac{x^k}{(1-x)^{k+1}}=\sum_{m\ge0}\binom{m+k}{k}x^{m+k}=\sum_{n\ge k}\binom nk x^n,
\]

所以 $[x^n]\frac{x^k}{(1-x)^{k+1}}=\binom nk$。

\subsection{运算表：生成函数怎么"翻译"数列操作}


设 $A(x)=\sum a_nx^n$、$B(x)=\sum b_nx^n$。下面是最常用的一组对应关系；每条都可以按系数逐项验证。

\textbf{线性组合。} $\{\alpha a_n+\beta b_n\}\longleftrightarrow \alpha A(x)+\beta B(x)$。

\textbf{右移。} $\{0,a_0,a_1,a_2,\dots\}\longleftrightarrow xA(x)$；更一般地，$\{0,\dots,0,a_0,a_1,\dots\}$（前面补 $m$ 个零）$\longleftrightarrow x^mA(x)$。

\textbf{左移。} 左移一格：

\[
\{a_1,a_2,a_3,\dots\}\longleftrightarrow \frac{A(x)-A(0)}{x}.
\]

若要从 $a_h$ 开始：

\[
\{a_h,a_{h+1},\dots\}\longleftrightarrow \frac{A(x)-a_0-a_1x-\cdots-a_{h-1}x^{h-1}}{x^h}.
\]

\textbf{数乘指数。} $\{\alpha^na_n\}\longleftrightarrow A(\alpha x)$。

\textbf{偶数项 / 奇数项。} 偶数下标：

\[
\{a_0,0,a_2,0,\dots\}\longleftrightarrow \frac{A(x)+A(-x)}{2};
\]

奇数下标：

\[
\{0,a_1,0,a_3,\dots\}\longleftrightarrow \frac{A(x)-A(-x)}{2}.
\]

\textbf{乘积（卷积）。} 设 $c_n=\sum_{i=0}^{n}a_ib_{n-i}$，则

\[
\{c_n\}\longleftrightarrow A(x)B(x).
\]

按系数展开即可：$[x^n]A(x)B(x)=\sum_{i=0}^na_ib_{n-i}$。

\textbf{求导。} $A'(x)=\sum_{n\ge1}na_nx^{n-1}=\sum_{n\ge0}(n+1)a_{n+1}x^n$。于是 $\{na_n\}\longleftrightarrow xA'(x)$。

\textbf{积分。} $\displaystyle\int_0^xA(t)\,dt=\sum_{n\ge0}\frac{a_n}{n+1}x^{n+1}$，即 $\left\{\frac{a_{n-1}}{n}\right\}_{n\ge1}\longleftrightarrow \displaystyle\int_0^xA(t)\,dt$。

\textbf{部分和。} 设 $s_n=a_0+a_1+\cdots+a_n$，则

\[
\{s_n\}\longleftrightarrow \frac{A(x)}{1-x},
\]

因为乘以 $\frac1{1-x}=1+x+x^2+\cdots$ 正是"对系数做部分和"。

这些运算合起来，就能把很多关于数列的递推/求和问题变成关于 $A(x)$ 的代数方程。

\subsection{例一：多米诺覆盖与斐波那契数}


\textbf{问题.} 用 $2\times1$ 的多米诺骨牌覆盖 $2\times n$ 的区域，有多少种覆盖方式？记这个数为 $f_n$，约定 $f_0=1$（什么都不放）。

\textbf{直接递推.} 看最左边：

\begin{itemize}
\item 若第一块是\textbf{竖}的，它占满第一列，剩下 $2\times(n-1)$ 的区域，有 $f_{n-1}$ 种；
\item 若第一块是\textbf{横}的，它必须与另一块横骨牌上下拼成一个 $2\times2$ 的方块，占掉前两列，剩下 $2\times(n-2)$，有 $f_{n-2}$ 种。
\end{itemize}


所以 $f_n=f_{n-1}+f_{n-2}$，初值 $f_0=f_1=1$。前几项是 $1,1,2,3,5,8,\dots$。

\textbf{生成函数解法.} 设 $F(x)=\sum_{n\ge0}f_nx^n$。把上面的分类"翻译"成生成函数：

\begin{itemize}
\item 空覆盖：贡献 $1$；
\item 第一块竖：贡献 $xF(x)$（占一列，用 $x$ 记一列）；
\item 第一块横（一对）：贡献 $x^2F(x)$（占两列）。
\end{itemize}


于是

\[
F(x)=1+xF(x)+x^2F(x)
\quad\Longrightarrow\quad
F(x)=\frac{1}{1-x-x^2}.
\]

这个式子把"递推关系"变成了"一个分式"；接下来只要把分式展开就能读出任意 $f_n$（下一节用部分分式给出通项）。

\textbf{两种骨牌的版本.} 若竖骨牌有 $x$ 种选择、横的一对有 $y$ 种选择，同样的分解给出

\[
F(x,y)=\frac{1}{1-x-y}=\sum_{m,n\ge0}\binom{m+n}{m}x^my^n.
\]

组合意义：一个覆盖由 $m$ 个竖块与 $n$ 个横对组成，共 $m+n$ 个"块"，选出哪 $m$ 个位置放竖块有 $\binom{m+n}{m}$ 种。

\subsection{例二：找零问题}


\textbf{问题.} 有 4 张 1 元硬币、无限多张 5 元纸币，凑出 $n$ 元有多少种方式？

分别写出两部分的生成函数：

\begin{itemize}
\item 4 张 1 元：最多用 4 张，方式数为 $1,1,1,1,1,0,0,\dots$，生成函数
\end{itemize}


\[
1+x+x^2+x^3+x^4=\frac{1-x^5}{1-x};
\]

\begin{itemize}
\item 无限 5 元：方式数为 $1,0,0,0,0,1,0,\dots$（0 元、5 元、10 元……），生成函数
\end{itemize}


\[
1+x^5+x^{10}+\cdots=\frac{1}{1-x^5}.
\]

（$n=0$ 时什么都不找，算 $1$ 种；这与两个生成函数的常数项都是 $1$ 一致。）两部分可以独立选择，所以"总方式数"的生成函数是两者\textbf{相乘}：

\[
\frac{1-x^5}{1-x}\cdot\frac{1}{1-x^5}=\frac{1}{1-x}.
\]

也就是说，凑出 $n$ 元恰好只有 $1$ 种方式——因为 1 元的张数只能是 $0,1,2,3,4$，而 $n$ 除以 5 的余数唯一确定了要用几张 1 元，剩下的用 5 元补足。这个例子说明：\textbf{生成函数的乘法对应"把 n 分成两部分、分别计数再求和"}；有时乘积会化简，从而直接看出答案。

\textbf{一般硬币问题.} 若有面值为 $w_1,\dots,w_k$ 的硬币各无限多张，凑出 $n$ 元的方式数的生成函数是

\[
\prod_{i=1}^{k}\frac{1}{1-x^{w_i}}.
\]

若第 $i$ 种硬币只有 $m_i$ 张，则对应因子换成 $\dfrac{1-x^{(m_i+1)w_i}}{1-x^{w_i}}$。整数分拆（把 $n$ 写成正整数之和，不计顺序）的生成函数就是

\[
\prod_{w\ge1}\frac{1}{1-x^w}.
\]

\subsection{本节的提醒}


\begin{itemize}
\item \textbf{生成函数是形式幂级数。} 写 $\frac{1}{1-x}$ 时不需要 $|x|<1$；只要按 $1+x+x^2+\cdots$ 理解即可。只有在做分析（如估计系数增长）时才需要收敛性。
\item \textbf{$x$ 记录什么要统一。} 在多米诺问题里 $x$ 记"一列的宽度"；若改成记"一块骨牌"，同一个问题的生成函数会变（通常变成 $x$ 与 $x^2$ 的多项式/分式）。定义清楚后，方程才不会列错。
\item \textbf{卷积的方向。} $[x^n]A(x)B(x)=\sum_{i=0}^na_ib_{n-i}$；若两边的"大小"含义不同（例如一边记列数、一边记块数），乘积的系数含义要先解释清楚。
\item \textbf{找零例子中"相乘"的前提是两部分独立。} 如果硬币之间有依赖（例如总共只能用 5 枚），就不能简单相乘，要引入新变量或改用别的模型。
\end{itemize}

\section{有理生成函数、部分分式与线性递推}


上一节把多米诺覆盖化成了 $F(x)=\dfrac{1}{1-x-x^2}$。这一节说明：\textbf{任何有理生成函数都能拆成若干 $\dfrac{1}{(1-\alpha x)^k}$ 之和}，从而把数列的通项写出来；反过来，\textbf{常系数线性递推}的生成函数一定是有理函数。这就是"用生成函数解递推"的完整流程。

\textbf{本节用到的记号：} $A(x),B(x),P(x),Q(x)$ 为多项式或形式幂级数；$\alpha,\beta$ 为常数；$[x^n]$ 为系数提取；$\binom nk$ 为二项式系数。

\subsection{有理函数与部分分式}


\textbf{定义.} 若 $A(x)=\dfrac{P(x)}{Q(x)}$，其中 $P,Q$ 是多项式、$Q(0)\neq0$，则称 $A(x)$ 是\textbf{有理生成函数}（或有理形式幂级数）。例如 $\dfrac{1}{1-x-x^2}$、$\dfrac{1+x}{(1-x)^3}$ 都是。

\textbf{基本事实（部分分式分解）.} 设 $Q(x)\neq0$ 且 $Q(0)\neq0$。在复数域上，$Q$ 可以分解成一次因子的乘积。则存在多项式 $R(x)$ 与常数 $c_{i,j}$，使

\[
\frac{P(x)}{Q(x)}=R(x)+\sum_{i=1}^{r}\sum_{j=1}^{k_i}\frac{c_{i,j}}{(1-\alpha_i x)^j},
\]

其中 $\alpha_i^{-1}$ 是 $Q$ 的根（重数 $k_i$）。若 $P$ 的次数小于 $Q$，则 $R=0$。

\textbf{为什么有用.} 因为每个 $\dfrac{1}{(1-\alpha x)^j}$ 的系数都知道：

\[
\frac{1}{(1-\alpha x)^j}=\sum_{n\ge0}\binom{n+j-1}{j-1}\alpha^n x^n.
\]

于是把部分分式拆开、逐项取系数，就得到原数列的通项公式（指数函数与多项式的乘积之和）。多项式部分 $R(x)$ 只影响有限项，后面可以忽略。

\textbf{操作步骤.} 求 $[x^n]P(x)/Q(x)$：

1. 若 $\deg P\ge\deg Q$，先做多项式除法，把 $P/Q$ 写成 $R+P_1/Q$，其中 $\deg P_1<\deg Q$；
2. 把 $Q$ 分解成一次因子的幂；
3. 设部分分式，比较系数（或代入特殊值）求出 $c_{i,j}$；
4. 对每个 $\dfrac{c}{(1-\alpha x)^j}$ 用上面的公式取 $x^n$ 的系数，相加。

\subsection{例一：斐波那契数的通项（Binet 公式）}


设 $f_0=f_1=1$、$f_n=f_{n-1}+f_{n-2}$，生成函数

\[
F(x)=\frac{1}{1-x-x^2}.
\]

\textbf{分解分母.} 令 $\varphi=\dfrac{1+\sqrt5}{2}$、$\psi=\dfrac{1-\sqrt5}{2}$，则 $\varphi+\psi=1$、$\varphi\psi=-1$，于是

\[
1-x-x^2=(1-\varphi x)(1-\psi x).
\]

\textbf{设部分分式.}

\[
\frac{1}{(1-\varphi x)(1-\psi x)}=\frac{A}{1-\varphi x}+\frac{B}{1-\psi x}.
\]

通分得 $A(1-\psi x)+B(1-\varphi x)=1$，即

\[
(A+B)-(A\psi+B\varphi)x=1.
\]

比较常数项与一次项：$A+B=1$、$A\psi+B\varphi=0$。解得

\[
A=\frac{\varphi}{\sqrt5},\qquad B=-\frac{\psi}{\sqrt5}.
\]

（也可用"代入 $x=1/\varphi$、$x=1/\psi$"的留数法快速求。）

\textbf{取系数.}

\[
F(x)=\frac{1}{\sqrt5}\left(\frac{\varphi}{1-\varphi x}-\frac{\psi}{1-\psi x}\right)
=\frac{1}{\sqrt5}\sum_{n\ge0}(\varphi^{n+1}-\psi^{n+1})x^n.
\]

所以

\[
f_n=\frac{\varphi^{n+1}-\psi^{n+1}}{\sqrt5}.
\]

这就是\textbf{Binet 公式}。验证：$n=0$ 时 $(\varphi-\psi)/\sqrt5=1$；$n=1$ 时 $(\varphi^2-\psi^2)/\sqrt5=(\varphi+\psi)=1$，都对。

\textbf{注意.} 这里的 $\varphi,\psi$ 是分母 $1-x-x^2$ 的根 $x$ 的倒数：$x^2=x+1$ 的正根是 $\varphi$，所以因子写成 $(1-\varphi x)$。做题时先把分母整理成"常数项为 $1$"的形式，再找它的根的倒数，比较不容易混。

\subsection{例二：一般二阶线性递推}


设 $a_n=\alpha a_{n-1}+\beta a_{n-2}$（$n\ge2$），初值 $a_0,a_1$ 给定。设 $A(x)=\sum a_nx^n$。由递推，

\[
A(x)-a_0-a_1x=\alpha x(A(x)-a_0)+\beta x^2A(x).
\]

整理：

\[
A(x)=\frac{a_0+(a_1-\alpha a_0)x}{1-\alpha x-\beta x^2}.
\]

分母是二次多项式。若 $1-\alpha x-\beta x^2=(1-\lambda_1x)(1-\lambda_2x)$ 有两个不同根，部分分式给出

\[
a_n=C_1\lambda_1^n+C_2\lambda_2^n;
\]

若重根 $\lambda$，则

\[
a_n=(C_1+C_2n)\lambda^n.
\]

这与第 12 讲用特征方程得到的形式完全一致；生成函数只是把"解递推"变成了"拆分式"。\textbf{重根时出现 $n\lambda^n$}，是因为部分分式里有 $\dfrac{1}{(1-\lambda x)^2}$ 这样的项，而它的系数是 $(n+1)\lambda^n$。

\textbf{一般结论（常系数线性递推与有理生成函数对应）.} 数列 $\{a_n\}$ 满足常系数线性递推

\[
a_n=c_1a_{n-1}+c_2a_{n-2}+\cdots+c_ra_{n-r}\qquad(n\ge r)
\]

当且仅当它的生成函数是有理函数

\[
A(x)=\frac{P(x)}{1-c_1x-c_2x^2-\cdots-c_rx^r},
\]

其中 $P$ 由初值决定、次数小于 $r$。这是因为递推两边乘以 $x^n$ 并对 $n\ge r$ 求和，右边恰好是 $A(x)$ 乘以 $1-c_1x-\cdots-c_rx^r$，左边只剩下有限个初值项。

\textbf{推论.} 常系数线性递推的解一定是"指数函数 × 多项式"的线性组合；反过来，任何有理生成函数对应的数列都满足某个常系数线性递推。

\subsection{例三：$\frac{1}{(1-x)^k}$ 与高阶部分分式}


由二项式定理（负指数版本），

\[
\frac{1}{(1-x)^k}=\sum_{n\ge0}\binom{n+k-1}{k-1}x^n.
\]

若分母有重根，例如

\[
\frac{1}{(1-2x)^3}=\sum_{n\ge0}\binom{n+2}{2}2^nx^n,
\]

系数是"多项式 × 指数"。这正是部分分式拆开后每一项的标准形式。

\textbf{综合例子.} 求 $[x^n]\dfrac{1+x}{(1-x)^3}$。先拆：

\[
\frac{1+x}{(1-x)^3}=\frac{2}{(1-x)^3}-\frac{1}{(1-x)^2}.
\]

因为 $1+x=2-(1-x)$。于是

\[
[x^n]\frac{1+x}{(1-x)^3}=2\binom{n+2}{2}-\binom{n+1}{1}
=2\cdot\frac{(n+2)(n+1)}{2}-(n+1)=(n+1)^2.
\]

所以 $\dfrac{1+x}{(1-x)^3}=\sum_{n\ge0}(n+1)^2x^n$。这个例子说明：\textbf{部分分式的系数不必真的解方程组，凑成已知形式往往更快}。

\subsection{本节的提醒}


\begin{itemize}
\item \textbf{先保证常数项为 $1$。} 部分分式的标准形式是 $\dfrac{1}{1-\alpha x}$；若分母常数项不是 $1$，先提出常数，或者改用 $(x-\alpha_i)$ 的形式并相应调整。
\item \textbf{部分分式在复数域上做。} 实系数多项式可能没有实根（如 $1+x^2$），但分解成一次因子需要复数；组合计数里允许用复数中间结果，最后实系数会抵消。
\item \textbf{有限项与无限项分开。} 若 $\deg P\ge\deg Q$，先做多项式除法；多项式部分只影响前有限项，不要漏掉。
\item \textbf{生成函数法与特征方程法等价。} 特征方程 $\lambda^2=\alpha\lambda+\beta$ 的根，就是分母 $1-\alpha x-\beta x^2$ 的根的倒数；两种写法不要混。
\end{itemize}

\section{卡特兰数的生成函数：从组合方程到通项}


第 12 讲已经用反射原理得到卡特兰数 $C_n=\dfrac{1}{n+1}\dbinom{2n}{n}$。这一节用生成函数\textbf{重新推出}这个公式：先由组合分解写出 $C(x)$ 满足的代数方程，再解方程、用广义二项式定理展开。这条路线是"生成函数解组合方程"的典型例子，也是下一节 $3\times n$ 多米诺覆盖方法的雏形。

\textbf{本节用到的记号：} $C_n$ 为第 $n$ 个卡特兰数（$C_0=1$：$n$ 对合法括号、Dyck 路、$(n+2)$ 边形三角剖分等）；$C(x)=\sum_{n\ge0}C_nx^n$；$\binom{\alpha}{k}$ 为广义二项式系数（$\alpha$ 可为任意实数）。

\subsection{组合分解给出代数方程}


\textbf{回顾卡特兰数的递推.} $C_0=1$，且对 $n\ge1$，

\[
C_n=\sum_{k=0}^{n-1}C_k\,C_{n-1-k}.
\]

\textbf{用生成函数翻译.} 设 $C(x)=\sum_{n\ge0}C_nx^n$。右边是一个卷积：$C(x)\cdot C(x)$ 的 $x^{n-1}$ 项系数正好是 $\sum_{k=0}^{n-1}C_kC_{n-1-k}=C_n$。因此对 $n\ge1$，$C_n=[x^{n-1}]C(x)^2$。乘以 $x^n$ 并对 $n\ge1$ 求和：

\[
C(x)-C_0=x\,C(x)^2.
\]

代入 $C_0=1$，得到\textbf{卡特兰生成函数方程}：

\[
C(x)=1+xC(x)^2.
\]

\textbf{组合意义.} 一个卡特兰对象（比如合法括号序列）可以唯一分解为

\[
\text{（一对括号包住的内部）}\ \text{后接}\ \text{（剩余部分）},
\]

内部是任意卡特兰对象（权 $C(x)$），外层括号与"后接"各贡献一个 $x$ 的"大小"，所以总生成函数满足 $C=1+xC^2$。这里的"$1$"对应空对象。$3\times n$ 覆盖、二叉树、三角剖分都可以做类似的分解，只是分解方式不同、方程不同。

\subsection{解二次方程}


把 $C=1+xC^2$ 整理成关于 $C$ 的二次方程：

\[
xC^2-C+1=0.
\]

由求根公式，

\[
C(x)=\frac{1\pm\sqrt{1-4x}}{2x}.
\]

\textbf{选哪个分支？} 生成函数在 $x=0$ 附近应等于 $C_0=1$。当 $x\to0$ 时，$\sqrt{1-4x}\to1$，分子 $1-\sqrt{1-4x}\to0$，与分母 $2x\to0$ 的比趋于有限值；而取"$+$"号时分子趋于 $2$，比值发散。因此取\textbf{减号}：

\[
C(x)=\frac{1-\sqrt{1-4x}}{2x}.
\]

这个式子本身已经是一个完整的答案：要算 $C_n$，只需把它展开。

\subsection{用广义二项式定理展开}


\textbf{广义二项式定理.} 对任意实数 $\alpha$ 与 $|x|$ 足够小（形式幂级数意义下则无需收敛），

\[
(1+x)^\alpha=\sum_{k\ge0}\binom{\alpha}{k}x^k,\qquad
\binom{\alpha}{k}=\frac{\alpha(\alpha-1)\cdots(\alpha-k+1)}{k!}.
\]

取 $\alpha=\frac12$、把 $x$ 换成 $-4x$：

\[
\sqrt{1-4x}=\sum_{k\ge0}\binom{1/2}{k}(-4x)^k.
\]

\textbf{求 $C_n$.} 由 $C(x)=\dfrac{1-\sqrt{1-4x}}{2x}$，

\[
C_n=[x^n]C(x)=\frac12[x^{n+1}]\left(1-\sqrt{1-4x}\right)
=-\frac12[x^{n+1}]\sqrt{1-4x}.
\]

（$[x^{n+1}]1=0$ 对 $n\ge0$。）代入展开式：

\[
C_n=-\frac12\binom{1/2}{n+1}(-4)^{n+1}.
\]

下面化简。先算 $\binom{1/2}{n+1}$：

\[
\binom{1/2}{n+1}
=\frac{\frac12\left(-\frac12\right)\left(-\frac32\right)\cdots\left(\frac12-n\right)}{(n+1)!}
=\frac{(-1)^n\,1\cdot3\cdot5\cdots(2n-1)}{2^{n+1}(n+1)!}.
\]

而 $1\cdot3\cdots(2n-1)=\dfrac{(2n)!}{2^nn!}$，所以

\[
\binom{1/2}{n+1}=\frac{(-1)^n(2n)!}{2^{2n+1}\,n!\,(n+1)!}.
\]

代回：

\[
C_n=-\frac12\cdot\frac{(-1)^n(2n)!}{2^{2n+1}n!(n+1)!}\cdot(-1)^{n+1}2^{2n+2}
=\frac{(2n)!}{n!(n+1)!}
=\frac{1}{n+1}\binom{2n}{n}.
\]

这就是卡特兰数的通项公式。\textbf{注意最后一步的因子}：$-\frac12\cdot(-1)^{2n+1}\cdot2^{2n+2}/2^{2n+1}=-\frac12\cdot(-1)\cdot2=1$，所以符号与 $2$ 的幂都恰好抵消——这是推导里最容易算错的地方。

\textbf{验证前几项.} 公式给出

\[
C_0,C_1,C_2,C_3,C_4,C_5=1,1,2,5,14,42,
\]

与递推一致。

\subsection{为什么这仍然是"普通"生成函数}


$C(x)=1+xC^2$ 里的乘法是普通卷积：把两个卡特兰对象按"大小相加"拼接。因此 $C(x)$ 是\textbf{普通生成函数}。若改用指数生成函数 $\hat C(x)=\sum C_nx^n/n!$，递推就不再做普通卷积，方程也不会这么简单（下一份讲解稿会专门对比这两种生成函数）。

\textbf{方法总结（组合方程 → 通项）.}

1. 找到组合对象的递归分解（"第一个块"/"最外层"/"根"等），写出 $C_n$ 与更小 $C_k$ 的关系；
2. 设生成函数 $C(x)$，把递推翻译成含 $C(x)$ 的代数方程（乘积对应卷积、移位对应乘 $x$）；
3. 解方程（二次方程用求根公式，注意用初值选分支）；
4. 用广义二项式定理/部分分式展开，读出系数。

这套流程对 $3\times n$ 覆盖、二叉树、平面树等都适用，只是第 1 步的分解不同。

\subsection{本节的提醒}


\begin{itemize}
\item \textbf{分支选择用初值，不用数值大小。} 选 $\frac{1-\sqrt{1-4x}}{2x}$ 是因为 $C(0)=1$；不要因为"$-\sqrt{}$ 看起来小"就选它。
\item \textbf{广义二项式系数要算对符号。} $\binom{1/2}{k}$ 的符号是 $(-1)^{k-1}$（$k\ge1$），与 $(-4)^k$ 的符号相乘后要仔细化简。
\item \textbf{$C=1+xC^2$ 中的 $x$ 记"一对括号"（或一个内部结点）的大小}；若换一个变量记法，方程形式会变，但方法不变。
\item \textbf{不要把它和"指数生成函数"的方程混淆。} 指数生成函数下，乘法的系数会多出二项式系数；直接套 $C=1+xC^2$ 会得到错误的数列。
\item \textbf{手写稿底部的方程字迹像 $x^2C^2$}，按上面的卷积推导应为 $xC^2$（一个 $x$ 记录"一对括号/一个内部结点"）；整理版统一写 $C=1+xC^2$。
\end{itemize}

\section{$3\times n$ 多米诺覆盖：辅助状态与生成函数方程组}


上一节的多米诺问题只涉及高度 $2$。若把区域高度改成 $3$，直接写一个递推就没那么容易了：第一列可能被"竖骨牌 + 横骨牌"混合填充，边界会留下缺口。课堂的处理办法是\textbf{引入辅助状态}，把"缺口"也当成一个生成函数，然后解方程组。这一节把这个过程完整写清楚。

\textbf{本节用到的记号：} $3\times n$ 表示高 $3$、宽 $n$ 的矩形；骨牌是 $2\times1$ 的多米诺；$x$ 记录"一列的宽度"；$C(x),V(x),W(x),Z(x)$ 是下面定义的四种轮廓状态的生成函数。

\subsection{问题与困难}


用 $2\times1$ 骨牌覆盖 $3\times n$ 的区域。因为总面积是 $3n$，必须为偶数，所以 $n$ 必须是\textbf{偶数}；$n=2$ 时有 $3$ 种覆盖，$n=4$ 时有 $11$ 种，$n=6$ 时有 $41$ 种。

困难在于：高度 $3$ 是奇数，第一列不可能只由竖骨牌填满，也不可能只由横骨牌填满（横骨牌一次占两列）；不同的填法会在边界留下不同的"缺口"，而缺口的状态会影响后面的填法。

\textbf{关于手写稿的方程.} 课堂手写稿用 V、Λ 两个"缺口状态"，但字面方程并不完整/不自洽：若照字面整理会得到 $C(x)=\dfrac{1-x^2}{1-2x^2-x^4}$，对应序列 $1,3,7,17,\dots$，与真实覆盖数 $1,3,11,41,153,\dots$ 不符。下面给出整理版重新推导的完整状态方程（C/V/W/Z 四种轮廓），解出的结果与直接枚举一致。

\subsection{四种轮廓状态}


从左往右一列一列地填。当已经填完前 $k$ 列、准备填第 $k+1$ 列时，有些骨牌可能从左边\textbf{伸进}第 $k+1$ 列。按第 $k+1$ 列已被占用的格子集合（称为\textbf{轮廓}）分类。对高度 $3$，可能的轮廓（忽略上下对称）只有四种：

\begin{itemize}
\item \textbf{C}：第 $k+1$ 列没有被占用的格子（干净边界）；
\item \textbf{V}：第 $k+1$ 列只有一个角上的格子被占用（"缺一个角"）；
\item \textbf{W}：第 $k+1$ 列有两个相邻的格子被占用（"缺两个格"）；
\item \textbf{Z}：第 $k+1$ 列的三个格子都已被占用（整列被横骨牌占满）。
\end{itemize}


记 $C(x),V(x),W(x),Z(x)$ 分别为从干净边界出发、经过若干列后到达相应轮廓、并且已经填好的部分的方式数的生成函数；$C(x)$ 中 $x^n$ 的系数就是 $3\times n$ 的完整覆盖数 $a_n$。

\textbf{为什么只有这四种？} 横骨牌只能占同一行的相邻两列，所以伸进下一列的格子必然是若干"整行"的集合；而竖骨牌在同一列里占相邻两行，会改变轮廓的形状。对高度 $3$，从干净边界出发能到达的轮廓只有上述四种（上下的对称算同一种，计数时要注意有 $2$ 种对称位置）。

\subsection{状态转移方程}


下面写出每种状态"填一列"后能变成什么。\textbf{每条转移乘一个 $x$，表示宽度增加一列。}

\textbf{从 C 出发.} 第 $k+1$ 列完全空，有三种填法：

1. 在两行放一个竖骨牌、剩下一行放一个横骨牌：竖骨牌可以在上两行或下两行，共 $2$ 种，各留下"一个角被占"的轮廓，贡献 $2xV(x)$；
2. 三行都放横骨牌：三个格子都被占，留下 Z 轮廓，贡献 $xZ(x)$。

再加上"$n=0$ 时什么都不放"的 $1$，得到

\[
C(x)=1+2xV(x)+xZ(x).
\]

\textbf{从 V 出发.} 假设被占的是上方（或下方）那个角上的格子。当前列还剩两个相邻格子：

1. 用一个竖骨牌填满这两个格子：回到干净边界，贡献 $xC(x)$；
2. 在两行各放一个横骨牌：伸出到下一列，留下"两个相邻格子被占"的轮廓 W，贡献 $xW(x)$。

所以

\[
V(x)=x\bigl(C(x)+W(x)\bigr).
\]

\textbf{从 W 出发.} 当前列还剩一个格子，只能放一个横骨牌伸到下一列，留下"一个角被占"的轮廓 V：

\[
W(x)=xV(x).
\]

\textbf{从 Z 出发.} 当前列已经被占满，不能再放任何骨牌，直接进入下一列，回到干净边界：

\[
Z(x)=xC(x).
\]

\subsection{解方程组}


把 $Z=xC$、$W=xV$ 代入：

\[
V=x(C+W)=xC+x^2V
\quad\Longrightarrow\quad
V=\frac{xC}{1-x^2}.
\]

再代入第一个方程：

\[
C=1+2xV+xZ
=1+\frac{2x^2C}{1-x^2}+x^2C.
\]

两边乘以 $1-x^2$：

\[
C(1-x^2)=(1-x^2)+2x^2C+x^2C(1-x^2).
\]

把含 $C$ 的项移到左边：

\[
C\bigl[(1-x^2)^2-2x^2\bigr]=1-x^2.
\]

展开 $(1-x^2)^2-2x^2=1-4x^2+x^4$，所以

\[
C(x)=\frac{1-x^2}{1-4x^2+x^4}.
\]

\subsection{验证与推论}


\textbf{展开前几项.} 用长除法或递推，

\[
C(x)=\frac{1-x^2}{1-4x^2+x^4}
=1+3x^2+11x^4+41x^6+153x^8+\cdots
\]

与直接数出的 $a_2=3$、$a_4=11$、$a_6=41$、$a_8=153$ 一致；奇数 $n$ 的系数全为 $0$，因为面积 $3n$ 必须是偶数。

\textbf{递推.} 由 $(1-4x^2+x^4)C(x)=1-x^2$，比较 $x^n$ 的系数：

\begin{itemize}
\item $n=0$：$a_0=1$；
\item $n=2$：$a_2-4a_0=-1$，即 $a_2=3$；
\item $n\ge4$：$a_n-4a_{n-2}+a_{n-4}=0$，即
\end{itemize}


\[
a_n=4a_{n-2}-a_{n-4}\qquad(n\ge4).
\]

验证：$a_4=4\cdot3-1=11$；$a_6=4\cdot11-3=41$；$a_8=4\cdot41-11=153$。

\textbf{这个例子的意义.} 它展示了生成函数方法的一般套路：

1. 把"填不完的边界"也当成状态，引入辅助生成函数；
2. 按"最后一列/第一列"的所有填法写转移方程；
3. 解线性方程组（这里是含 $C,V,W,Z$ 的一次方程组），得到主生成函数的有理表达式；
4. 由有理表达式读出递推与数值。

对更高的矩形（如 $4\times n$、$5\times n$）或别的骨牌形状，只要轮廓状态有限，都可以照这个办法列方程；状态一多，手算会变繁，但思路不变。计算机算法里的"轮廓线 DP/转移矩阵"就是这套生成函数方程的系数版本。

\subsection{本节的提醒}


\begin{itemize}
\item \textbf{$x$ 记的是列宽，不是骨牌数。} 因此 $C(x)$ 只有偶次项；若改成"每块骨牌乘 $x$"，表达式会变，但系数含义也随之改变。
\item \textbf{对称状态要乘 $2$。} "缺一个角"有上、下两种位置，"缺两个格"也有上、下两种位置；在方程里已经用系数 $2$ 处理，写方程时不要漏。
\item \textbf{状态要选"能写出转移"的。} 缺口状态不是"为了好看"随便加的，而是因为直接对完整矩形写递推会漏情况；把边界状态补齐后，每种填法恰好落入某一类，方程才不重不漏。
\item \textbf{解出的分母给出递推。} 有理生成函数的分母对应线性递推；这和第 2 节"常系数线性递推与有理生成函数的对应"是同一件事。
\end{itemize}

\section{指数生成函数与带标号结构}


前几节的生成函数都是 $A(x)=\sum a_nx^n$，乘法对应普通卷积。这一节引入\textbf{指数生成函数}（exponential generating function，EGF）

\[
\hat A(x)=\sum_{n\ge0}a_n\frac{x^n}{n!}.
\]

多除一个 $n!$，换来的是：\textbf{乘法对应"二项卷积"}，而二项卷积正好描述"把 $n$ 个带标签的对象分成两部分"。因此 EGF 是数\textbf{带标号结构}（排列、映射、集合划分、标号树等）的专用工具。

\textbf{本节用到的记号：} $a_n$ 为某个组合对象在 $n$ 元集合上的个数；$\hat A(x)=\sum a_nx^n/n!$；$[x^n]$ 为系数提取；$\binom ni$ 为二项式系数；$S(n,k)$ 为第二类 Stirling 数；$D_n$ 为错排数。

\subsection{定义与第一个例子}


\textbf{定义.} 数列 $\{a_n\}$ 的\textbf{指数生成函数}是

\[
\hat A(x)=\sum_{n\ge0}a_n\frac{x^n}{n!}.
\]

注意 $a_n=n![x^n]\hat A(x)$。例如：

\begin{itemize}
\item $a_n=1$（所有 $n$ 只有一种结构）：$\hat A(x)=\sum x^n/n!=e^x$；
\item $a_n=n!$（$n$ 元排列）：$\hat A(x)=\sum n!\,x^n/n!=\sum x^n=\dfrac{1}{1-x}$；
\item $a_n=\binom nk$：$\hat A(x)=\sum\binom nk\frac{x^n}{n!}=e^x\dfrac{x^k}{k!}$（因为 $\binom nk/n!=1/(k!(n-k)!)$）。
\end{itemize}


对比：同样一个数列，用 OGF 和 EGF 得到的闭式往往完全不同。例如 $a_n=n!$ 的 OGF 是 $\sum n!x^n$（没有简单的初等闭式），但 EGF 是 $\frac{1}{1-x}$，非常好用。这说明\textbf{选对生成函数类型}很重要。

\subsection{乘法：二项卷积}


\textbf{定理（EGF 乘法）.} 设 $\hat A(x)=\sum a_nx^n/n!$、$\hat B(x)=\sum b_nx^n/n!$，则

\[
n!\,[x^n]\hat A(x)\hat B(x)=\sum_{i=0}^{n}\binom ni a_i b_{n-i}.
\]

\textbf{证明.} 直接算：

\[
\hat A\hat B
=\sum_{i,j\ge0}\frac{a_i}{i!}\frac{b_j}{j!}x^{i+j}
=\sum_{n\ge0}\left(\sum_{i=0}^{n}\frac{a_ib_{n-i}}{i!(n-i)!}\right)x^n.
\]

乘以 $n!$ 并把 $n!/(i!(n-i)!)$ 写成 $\binom ni$，即得。$\square$

\textbf{组合意义（关键）.} 若 $a_i$ 是"在 $i$ 个带标签的元素上造一个 $\mathcal A$ 结构"的个数，$b_j$ 同理，那么

\[
\binom ni a_ib_{n-i}
\]

数的是：\textbf{从 $n$ 个标签中选出 $i$ 个分给 $\mathcal A$ 结构，剩下 $n-i$ 个分给 $\mathcal B$ 结构，再分别在各自标签集上造结构}。对所有 $i$ 求和，就是"把 $n$ 个标签分成两部分、分别造 $\mathcal A$ 与 $\mathcal B$ 结构"的总数。因此

\[
\hat A\hat B \ \text{的系数}\times n! \ =\ \text{把标签分成两部分分别造结构的方案数}.
\]

这就是 EGF 乘法的含义：\textbf{先分标签，再做结构}。OGF 的普通卷积不区分标签，因此不能用来数带标号对象。

\subsection{平移、求导与 $\exp$}


\textbf{求导（平移）.} 若 $\hat A(x)=\sum a_nx^n/n!$，则

\[
\hat A'(x)=\sum_{n\ge1}a_n\frac{x^{n-1}}{(n-1)!}=\sum_{m\ge0}a_{m+1}\frac{x^m}{m!},
\]

即\textbf{求导等于数列左移一格}。这是 EGF 比 OGF 方便的地方之一：OGF 求导会带出因子 $n$，而 EGF 求导是干净的平移。

\textbf{多项式的算子写法.} 若 $P$ 是多项式，则 $\{P(n)a_n\}$ 的 EGF 是 $P\!\left(x\dfrac{d}{dx}\right)\hat A(x)$，因为 $x\frac{d}{dx}$ 把 $a_n$ 变成 $n a_n$。

\textbf{指数公式（集合构造）.} 设 $\hat G(x)$ 是"连通/不可分结构"的 EGF。把若干个这样的结构\textbf{无序地}拼成一个整体（称为集合构造，SET），其 EGF 为

\[
e^{\hat G(x)}=1+\hat G+\frac{\hat G^2}{2!}+\frac{\hat G^3}{3!}+\cdots
\]

\textbf{解释.} $\hat G^k$ 数的是"把标签分成 $k$ 个有序部分、每部分造一个 $\mathcal G$ 结构"（因为每次乘法都先分标签）；但集合里的 $k$ 个部分\textbf{没有顺序}，所以要除以 $k!$。对 $k$ 求和即得 $e^{\hat G}$。类似地，\textbf{有序序列}（SEQ）的 EGF 是

\[
\frac{1}{1-\hat G}=1+\hat G+\hat G^2+\cdots
\]

（这里不除以 $k!$，因为序列有顺序）。

\subsection{例一：排列、错排、集合划分}


\textbf{排列.} 一个 $n$ 元排列可以看成"若干个轮换（cycle）的无序集合"。$k$ 元轮换的个数是 $(k-1)!$，所以轮换的 EGF 是

\[
\sum_{k\ge1}(k-1)!\frac{x^k}{k!}=\sum_{k\ge1}\frac{x^k}{k}=-\log(1-x).
\]

由集合构造公式，排列的 EGF 是

\[
\exp\bigl(-\log(1-x)\bigr)=\frac{1}{1-x},
\]

其 $x^n$ 系数乘以 $n!$ 得 $n!$，正是排列数。验证了公式。

\textbf{错排.} 错排是"没有长度 $1$ 轮换"的排列，即只由长度 $\ge2$ 的轮换组成的集合。轮换 EGF 去掉长度 $1$ 的项 $x$，得 $-\log(1-x)-x$，于是错排 EGF 为

\[
\exp\bigl(-\log(1-x)-x\bigr)=\frac{e^{-x}}{1-x}.
\]

展开：$e^{-x}/(1-x)=\sum_n\left(\sum_{k=0}^n(-1)^k/k!\right)x^n$，乘以 $n!$ 得

\[
D_n=n!\sum_{k=0}^{n}\frac{(-1)^k}{k!},
\]

与第 12 讲容斥得到的公式一致。

\textbf{集合划分与 Bell 数.} 非空块（一个块就是一个非空集合）的 EGF 是 $e^x-1$（$n$ 元集合只有一种结构，$n\ge1$）。集合划分是"非空块的无序集合"，所以 Bell 数的 EGF 是

\[
\exp(e^x-1).
\]

前几项：$B_0,B_1,B_2,\dots=1,1,2,5,15,52,203,\dots$。

\textbf{满射.} 从 $n$ 元集到 $m$ 元集的满射：把 $n$ 个标签分成 $m$ 个非空有序块（每块是某个原像），每块的 EGF 是 $e^x-1$，所以满射数的 EGF 是 $(e^x-1)^m$。其 $x^n$ 系数乘以 $n!$ 等于 $\sum_i(-1)^i\binom mi(m-i)^n=m!S(n,m)$，与容斥公式一致。

\subsection{例二：教室与"集合"的直观解释}


课堂上的例子：有 $n$ 个\textbf{不同的}学生，要分到若干间\textbf{教室}；若一间教室里有 $j$ 个学生，这 $j$ 个学生在教室内部有 $g_j$ 种安排方式（约定 $g_0=0$，空教室不计）。设

\[
\hat G(x)=\sum_{j\ge0}g_j\frac{x^j}{j!}.
\]

\begin{itemize}
\item 若教室是\textbf{有编号}的，且恰好用 $k$ 间：先把 $n$ 个标签分成 $k$ 个有序部分（每部分对应一间教室），每部分内部造 $\mathcal G$ 结构，总数正是 $n![x^n]\hat G^k$；
\item 若教室\textbf{没有编号}：$k$ 间教室的排列不区分，所以要除以 $k!$，总数是 $n![x^n]\dfrac{\hat G^k}{k!}$；
\item 若教室数量不限：对 $k$ 求和，得到 $n![x^n]e^{\hat G(x)}$。
\end{itemize}


这就是 $e^{\hat G}$ 的"集合构造"含义：\textbf{把带标签的对象分成若干个无序的块，每块内部是一个 $\mathcal G$ 结构}。

\subsection{例三：有根树（每个结点最多 3 个无序孩子）}


设 $t_n$ 是 $n$ 个结点的\textbf{带标号有根树}个数，每个结点最多有 $3$ 个孩子，孩子之间\textbf{不区分顺序}。设 $\hat T(x)=\sum t_nx^n/n!$。

一棵树 = 根 + 一个"由 $0,1,2,3$ 棵子树组成的无序集合"：
\begin{itemize}
\item $0$ 棵子树：$1$；
\item $1$ 棵：$\hat T$；
\item $2$ 棵：两棵子树无序，EGF 为 $\hat T^2/2!$；
\item $3$ 棵：$\hat T^3/3!$。
\end{itemize}


再乘上根自己的标签 $x$，得

\[
\hat T(x)=x\left(1+\hat T+\frac{\hat T^2}{2!}+\frac{\hat T^3}{3!}\right).
\]

展开前几项（用 $t_n=n![x^n]\hat T$）：

\[
t_1,t_2,t_3,t_4,t_5,t_6=1,2,9,64,620,7620.
\]

作为对比，若孩子数量不限，方程是 $\hat T=xe^{\hat T}$，其系数为 $t_n=n^{n-1}$（Cayley 公式的带标号有根树版本）：$1,2,9,64,625,7776,\dots$。前 $4$ 项相同（因为 $n\le4$ 时不可能有 $4$ 个孩子），从 $n=5$ 开始分叉（$620$ 对 $625$）。这个对比说明：\textbf{"最多 3 个孩子"这个限制只影响 $n\ge5$ 的项}，EGF 方程精确地反映了这一点。

\textbf{注意标号与无标号的区别.} 上面数的是\textbf{带标号}的有根树（$n$ 个结点带 $1,\dots,n$ 的标签），所以用 EGF。若问的是\textbf{无标号}的形状（同构的树算一种），则要换成 OGF，且\textbf{不能}直接写 $T=x(1+T+T^2+T^3)$——那个方程数的是"孩子有顺序"的平面树，前几项是 $1,1,2,5,13,36,\dots$。

无标号、孩子无序且最多 $3$ 个的 OGF 要用多重集（Polya）构造：

\[
T(x)=x\left[1+T(x)+\frac{T(x)^2+T(x^2)}{2}+\frac{T(x)^3+3T(x)T(x^2)+2T(x^3)}{6}\right],
\]

前几项是 $1,1,2,4,8,17,39,\dots$（$n=1,2,\dots$）。另外，$T=x(1+T+T^2+T^3)$ 的系数 $1,1,2,5,13,\dots$ 与"无限制无标号有根树"的 $1,1,2,4,9,20,\dots$（A000081）也不是同一列数：三者（标号 EGF / 有序孩子 OGF / 无标号无序孩子 OGF）必须分清。课堂转写这段没有讲完，整理版按"带标号树"叙述，并在此注明无标号版本的区别。

\subsection{本节与 OGF 的对照（先给结论）}


\begin{itemize}
\item \textbf{OGF 的乘法是普通卷积}：$\sum a_ib_{n-i}$，对应"把总大小 $n$ 分成 $i$ 和 $n-i$，两部分不区分标签"；
\item \textbf{EGF 的乘法是二项卷积}：$\sum\binom ni a_ib_{n-i}$，对应"先把 $n$ 个标签分成两部分，再分别造结构"；
\item 因此：数\textbf{带标号}对象（排列、映射、集合划分、标号树、满射）用 EGF；数\textbf{无标号}对象（分拆、整数解、格路、括号、Catalan、找零）用 OGF；
\item 结构是"无序集合"时，EGF 用 $\exp$；结构是"有序列"时，EGF 用 $1/(1-\hat G)$；相应的 OGF 运算则是普通乘积/复合。
\end{itemize}


更系统的对比、判别准则与常见错误，见附加讲解稿《普通生成函数 vs 指数生成函数》。

\subsection{本节的提醒}


\begin{itemize}
\item \textbf{EGF 的系数要乘回 $n!$。} 算出的 $\hat A$ 的 $x^n$ 系数是 $a_n/n!$；答案要的是 $a_n$，别忘了乘以 $n!$。
\item \textbf{$\exp$ 只对 EGF 成立。} 对 OGF 写 $e^{A(x)}$ 通常没有组合意义；Catalan 的 OGF 满足 $C=1+xC^2$，不能写成 $1/(1-C)$ 之类。
\item \textbf{"标签"要真的可区分。} 若对象没有标签（如相同的球、整数分拆），用 EGF 会多除 $n!$，答案会错。
\item \textbf{有序与无序。} 有序序列用 $1/(1-\hat G)$，无序集合用 $e^{\hat G}$；两者差一个 $k!$，对应"是否给块编号"。
\end{itemize}

\section{幂和与 Bernoulli 数：用 EGF 把求和变成乘法}


这一节是 EGF 的一个漂亮应用：求

\[
S_m(n)=1^m+2^m+\cdots+n^m
\]

的\textbf{关于 $n$ 的公式}（$m$ 固定）。直接对每个 $m$ 硬算很麻烦；若把 $S_m(n)$ 看成"以 $m$ 为下标"的数列，它的\textbf{指数生成函数}恰好可以写成两个已知函数的乘积，乘积一展开，公式就出来了。

\textbf{本节用到的记号：} $S_m(n)=\sum_{k=1}^{n}k^m$；$B_j$ 为 Bernoulli 数（本讲采用 $B_1=+\frac12$ 的约定，记作 $B_j^+$，见下）；$[x^n]$ 为系数提取；$\binom nk$ 为二项式系数。

\subsection{关键恒等式：幂和的 EGF}


对固定的 $n$，考虑以 $m$ 为下标的数列 $\{S_m(n)\}_{m\ge0}$。它的指数生成函数是

\[
\sum_{m\ge0}S_m(n)\frac{x^m}{m!}
=\sum_{m\ge0}\left(\sum_{k=1}^{n}k^m\right)\frac{x^m}{m!}
=\sum_{k=1}^{n}\sum_{m\ge0}\frac{(kx)^m}{m!}
=\sum_{k=1}^{n}e^{kx}.
\]

（第一步交换求和顺序；第二步用 $e^{kx}=\sum_m (kx)^m/m!$。）右边是等比级数：

\[
\sum_{k=1}^{n}e^{kx}=e^x\frac{e^{nx}-1}{e^x-1}
=\left(\frac{xe^x}{e^x-1}\right)\cdot\left(\frac{e^{nx}-1}{x}\right).
\]

这两个因子分别有干净的展开：

\[
\frac{xe^x}{e^x-1}=\sum_{j\ge0}B_j^+\frac{x^j}{j!},
\qquad
\frac{e^{nx}-1}{x}=\sum_{i\ge0}\frac{n^{i+1}x^i}{(i+1)!}.
\]

第一个定义的就是 \textbf{Bernoulli 数}（采用"加号约定"，见下）；第二个只是把 $e^{nx}-1$ 的展开除以 $x$。于是\textbf{幂和的 EGF 是两个已知级数的乘积}，这就是它好算的原因。

\subsection{Bernoulli 数的定义与计算}


\textbf{定义（加号约定）.} 定义 $B_j^+$ 为

\[
\frac{xe^x}{e^x-1}=\sum_{j\ge0}B_j^+\frac{x^j}{j!}.
\]

前几个值：

\[
B_0^+=1,\quad B_1^+=\frac12,\quad B_2^+=\frac16,\quad B_3^+=0,\quad B_4^+=-\frac1{30},\quad B_5^+=0,\quad B_6^+=\frac1{42}.
\]

（有的书用 $\dfrac{x}{e^x-1}=\sum B_j\frac{x^j}{j!}$，此时 $B_1=-\frac12$，其余 $B_j$ 与 $B_j^+$ 在 $j\ne1$ 时相同。两种约定只差 $B_1$ 的符号，公式里要配套使用。）

\textbf{递推.} 由

\[
\frac{xe^x}{e^x-1}\cdot\frac{e^x-1}{x}=e^x
\]

两边展开、比较 $x^m/m!$ 的系数，可得

\[
\sum_{j=0}^{m}\binom{m+1}{j}B_j^+=m+1\qquad(m\ge0).
\]

前三项验证：$m=0$：$B_0^+=1$；$m=1$：$B_0^++2B_1^+=2$，即 $B_1^+=\frac12$；$m=2$：$B_0^++3B_1^++3B_2^+=3$，即 $B_2^+=\frac16$。这个递推可以依次算出所有 $B_j^+$，且奇数下标 $j\ge3$ 的 Bernoulli 数为 $0$。

\subsection{幂和公式}


把上面的乘积按 $x^m/m!$ 取系数。设 $j+i=m$，则

\[
S_m(n)=\sum_{j=0}^{m}\frac{m!}{j!\,(i+1)!}B_j^+\,n^{i+1}
=\frac{1}{m+1}\sum_{j=0}^{m}\binom{m+1}{j}B_j^+\,n^{m+1-j}.
\]

\textbf{这就是 Bernoulli 幂和公式。}

\textbf{验证几个低次情形.}

\begin{itemize}
\item $m=0$：$S_0(n)=\dfrac{1}{1}\dbinom10B_0^+n=n$（共 $n$ 项，每项 $1$）。
\item $m=1$：$S_1(n)=\dfrac12\left(B_0^+n^2+2B_1^+n\right)=\dfrac{n^2+n}{2}=\dfrac{n(n+1)}2$。
\item $m=2$：$S_2(n)=\dfrac13\left(B_0^+n^3+3B_1^+n^2+3B_2^+n\right)=\dfrac13\left(n^3+\dfrac32n^2+\dfrac12n\right)=\dfrac{n(n+1)(2n+1)}6$。
\item $m=3$：$S_3(n)=\dfrac14\left(n^4+2n^3+n^2\right)=\left(\dfrac{n(n+1)}2\right)^2$。
\item $m=4$：$S_4(n)=\dfrac15\left(n^5+\dfrac52n^4+\dfrac53n^3-\dfrac16n\right)=\dfrac{n(n+1)(2n+1)(3n^2+3n-1)}{30}$。
\end{itemize}


这些都与常用的求和公式一致。

\subsection{为什么这比"背公式"省事}


对每个固定的 $m$，$S_m(n)$ 是关于 $n$ 的 $m+1$ 次多项式，且常数项为 $0$（因为 $S_m(0)=0$）。若用待定系数法，每个 $m$ 都要解一个 $m+2$ 元线性方程组；而用上面的公式，\textbf{只需要一套 Bernoulli 数 $B_0^+,B_1^+,\dots$}，代入即可。$m$ 增大时，Bernoulli 数可以递推得到，计算量是"一次算好、反复使用"。

课堂手写稿里把这写成

\[
S_m(n)=m!\sum_{j=0}^{m}\frac{B_j^+\,n^{m+1-j}}{j!\,(m+1-j)!},
\]

它与上面的公式只差把 $\frac{1}{m+1}\binom{m+1}{j}$ 拆成 $\frac{m!}{j!(m+1-j)!}$，是同一个式子。

\subsection{与 EGF 主题的联系}


这个例子把 EGF 的两个特点都用上了：

1. \textbf{EGF 的乘法是二项卷积}：两个级数相乘后，$x^m/m!$ 的系数自动带出 $\binom{m+?}{?}$ 的组合系数，正好对应公式里的二项式系数；
2. \textbf{指数函数的 Taylor 展开}把"离散求和 $\sum_{k=1}^n$"变成了"等比级数"，从而可以写成乘积。

如果一开始把 $S_m(n)$ 看成普通数列（以 $n$ 或 $m$ 为下标）用 OGF 处理，就很难得到这样干净的乘积结构；这正是 EGF 适合"带组合系数/带标号"问题的体现。

\subsection{本节的提醒}


\begin{itemize}
\item \textbf{Bernoulli 约定要统一。} 用 $B_1^+=\frac12$ 时公式里全是加号；若用 $B_1=-\frac12$，公式中 $j=1$ 的项要变号。不要两套约定混用。
\item \textbf{$S_m(n)$ 是 $n$ 的多项式，不是 $m$ 的多项式。} 公式对固定 $m$、任意 $n$ 成立；对调 $m,n$ 不成立。
\item \textbf{常数项为 $0$。} 因为 $n=0$ 时和为 $0$；用公式算出的多项式应满足这一点，可作为检查。
\item \textbf{$m=0$ 的约定。} 按 $0^0=1$（组合上通常约定）或直接数项数，$S_0(n)=n$；公式给出 $n$，一致。
\item \textbf{手写稿的端点有 off-by-one。} 手写稿把 $\sum_{k=1}^ne^{kx}$ 写成 $\dfrac{e^{nx}-1}{e^x-1}$ 一类的形式，漏掉了 $k=1$ 那项的 $e^x$ 因子；正确是 $\dfrac{e^{(n+1)x}-e^x}{e^x-1}=e^x\dfrac{e^{nx}-1}{e^x-1}$。整理版以本节的式子为准。
\end{itemize}


\section*{一页式总览}
\begin{itemize}
\item 普通生成函数：$A(x)=\sum_{n\ge0}a_nx^n$，$a_n=[x^n]A$；形式幂级数，不必考虑收敛。
\item OGF 基本表：$1/(1-x)$、$1/(1-\alpha x)$、$1/(1-x^p)$、$x^p/(1-x)^{p+1}\leftrightarrow\binom np$、$1/(1-x)^k\leftrightarrow\binom{n+k-1}{k-1}$。
\item OGF 运算：右移 $xA$、左移 $(A-A(0))/x$、$\alpha^na_n\leftrightarrow A(\alpha x)$、偶/奇子列 $(A(x)\pm A(-x))/2$、乘积 $=$ 卷积、$A'(x)=\sum(n+1)a_{n+1}x^n$、$xA'(x)=\sum na_nx^n$、部分和 $A/(1-x)$。
\item 多米诺/斐波那契：$F=1+xF+x^2F=\dfrac1{1-x-x^2}$；两种骨牌 $\dfrac1{1-x-y}=\sum\binom{m+n}{m}x^my^n$。
\item 找零：$\prod_i 1/(1-x^{w_i})$；4 张 1 元 + 无限 5 元 $\Rightarrow (1-x^5)/(1-x)\cdot 1/(1-x^5)=1/(1-x)$。
\item 部分分式：有理函数拆成 $\sum c_{ij}/(1-\alpha_ix)^j$，用 $1/(1-\alpha x)^j=\sum_n\binom{n+j-1}{j-1}\alpha^nx^n$ 取系数；斐波那契 $f_n=(\varphi^{n+1}-\psi^{n+1})/\sqrt5$。
\item Catalan：$C=1+xC^2\Rightarrow C=(1-\sqrt{1-4x})/(2x)$；广义二项式展开得 $C_n=\frac1{n+1}\binom{2n}{n}$。
\item $3\times n$ 覆盖：轮廓状态 C/V/W/Z，$C=1+2xV+xZ$、$V=x(C+W)$、$W=xV$、$Z=xC$；解得 $C(x)=\dfrac{1-x^2}{1-4x^2+x^4}$，$a_n=4a_{n-2}-a_{n-4}$（$n$ 为偶数）。
\item 指数生成函数：$\hat A(x)=\sum a_nx^n/n!$；乘法 $n![x^n]\hat A\hat B=\sum_i\binom ni a_ib_{n-i}$；求导=左移；$\exp(\hat G)$ 数"无序集合"，$1/(1-\hat G)$ 数"有序序列"。
\item EGF 例：$1\leftrightarrow e^x$、$n!\leftrightarrow 1/(1-x)$、错排 $\leftrightarrow e^{-x}/(1-x)$、Bell $\leftrightarrow e^{e^x-1}$、满射 $\leftrightarrow (e^x-1)^m$；$\le3$ 个无序孩子的有根树 $\hat T=x(1+\hat T+\hat T^2/2!+\hat T^3/3!)$。
\item 幂和：$\sum_{m\ge0}S_m(n)\frac{x^m}{m!}=e^x\dfrac{e^{nx}-1}{e^x-1}$。
\item Bernoulli：$B_j^+$ 由 $\dfrac{xe^x}{e^x-1}=\sum_jB_j^+x^j/j!$ 定义；$S_m(n)=\dfrac1{m+1}\sum_{j=0}^m\binom{m+1}{j}B_j^+n^{m+1-j}$。
\end{itemize}

\section*{参考资料}
\begin{itemize}
\item \textbf{[DS]} 邓小铁、王畅。《离散数学与结构》第 13 章：§13.1 用生成函数求解递归式、§13.2 一般生成函数、§13.3 指数生成函数。
\item \textbf{[讲义]} 课程讲义 \texttt{25-26-lecture\_notes.pdf} §9.9 生成函数（定义、基本表、运算、找零、部分分式、Fibonacci、Catalan）。
\item 原始材料：本地归档（手写原稿与录音转写），不纳入课程仓库。
\item 附加讲解稿：本目录 \texttt{explainer-ogf-vs-egf.pdf}。
\end{itemize}
\end{document}
